\documentclass[12pt, titlepage]{article}

\usepackage{amsmath, mathtools}

\usepackage[round]{natbib}
\usepackage{amsfonts}
\usepackage{amssymb}
\usepackage{graphicx}
\usepackage{colortbl}
\usepackage{xr}
\usepackage{hyperref}
\usepackage{longtable}
\usepackage{xfrac}
\usepackage{tabularx}
\usepackage{float}
\usepackage{siunitx}
\usepackage{booktabs}
\usepackage{multirow}
\usepackage[section]{placeins}
\usepackage{caption}
\usepackage{fullpage}

\hypersetup{
bookmarks=true,     % show bookmarks bar?
colorlinks=true,       % false: boxed links; true: colored links
linkcolor=red,          % color of internal links (change box color with linkbordercolor)
citecolor=blue,      % color of links to bibliography
filecolor=magenta,  % color of file links
urlcolor=cyan          % color of external links
}

\usepackage{array}

\externaldocument{../../SRS/SRS}

\input{../../Comments.text}
\input{../../Common.text}

\begin{document}

\title{Module Interface Specification for \progname{}}

\author{\authname}

\date{\today}

\maketitle

\pagenumbering{roman}

\section{Revision History}

\begin{tabularx}{\textwidth}{p{3cm}p{2cm}X}
\toprule {\bf Date} & {\bf Version} & {\bf Notes}\\
\midrule
Mar 25, 2026 & 1.1 & Post-presentation refinements based on professor's feedback\\
March 5, 2026 & 1.0 & Initial version\\
\bottomrule
\end{tabularx}

~\newpage

\section{Symbols, Abbreviations and Acronyms}

See SRS documentation at
\href{https://github.com/Daffodil404/yanyu_wu_cas741/blob/main/docs/SRS/SRS.tex}{SRS.tex}.

% \wss{Also add any additional symbols, abbreviations or acronyms}

\newpage

\tableofcontents

\newpage

\pagenumbering{arabic}

\section{Introduction}

The following document details the Module Interface Specifications for integrating NCRF into the existing TRF pipeline. This pipeline accepts MEG/EEG data for analyzing, and we now add estimators to distinguish different estimation methods, such as \texttt{ncrf} and \texttt{boosting}.
% \wss{Fill in your project name and description}

Complementary documents include the System Requirement Specifications
and Module Guide.  The full documentation and implementation can be
found at \url{https://github.com/Daffodil404/yanyu_wu_cas741}. The key
companion documents are
\href{https://github.com/Daffodil404/yanyu_wu_cas741/blob/main/docs/SRS/SRS.tex}{SRS}
and
\href{https://github.com/Daffodil404/yanyu_wu_cas741/blob/main/docs/Design/SoftArchitecture/MG.tex}{MG}.

\section{Notation}

% \wss{You should describe your notation.  You can use what is below as
%   a starting point.}

The structure of the MIS for modules comes from \citet{HoffmanAndStrooper1995},
with the addition that template modules have been adapted from
\cite{GhezziEtAl2003}.  The mathematical notation comes from Chapter 3 of
\citet{HoffmanAndStrooper1995}. 
%  For instance, the symbol := is used for a
% multiple assignment statement and conditional rules follow the form $(c_1
% \Rightarrow r_1 | c_2 \Rightarrow r_2 | ... | c_n \Rightarrow r_n )$.

The following table summarizes the primitive data types used by \progname. 

\begin{center}
\renewcommand{\arraystretch}{1.2}
\noindent 
\begin{tabular}{l l p{7.5cm}} 
\toprule 
\textbf{Data Type} & \textbf{Notation} & \textbf{Description}\\ 
\midrule
boolean & $\mathbb{B}$ & logical value: \texttt{true} or \texttt{false}\\
integer & $\mathbb{Z}$ & a whole number in $(-\infty, \infty)$\\
natural number & $\mathbb{N}$ & a counting number in $[1, \infty)$\\
real scalar & $\mathbb{R}$ & a real-valued scalar quantity\\
real vector & $\mathbb{R}^{n}$ & a one-dimensional real-valued time series\\
real matrix & $\mathbb{R}^{m \times n}$ & a two-dimensional real-valued array\\
real tensor & $\mathbb{R}^{m \times i \times \ell}$ & a three-dimensional real-valued array\\
EEG/MEG signal & $\mathbb{R}^{c \times t}$ & multichannel real-valued time series, where $c$ is channel/sensor count and $t$ is time index\\
string & str & text identifier such as subject, task, or estimator name\\
path & Path & filesystem location for input/output artifacts\\
\bottomrule
\end{tabular} 
\end{center}

\noindent
The MIS specification of \progname{} uses the primitive data types listed in
the table above: $\mathbb{B}$ (boolean), $\mathbb{Z}$ (integer),
$\mathbb{N}$ (natural number), $\mathbb{R}$ (real scalar),
$\mathbb{R}^{n}$ (real vector), $\mathbb{R}^{m \times n}$ (real matrix), and
$\mathbb{R}^{m \times i \times \ell}$ (real tensor), plus \texttt{str} and
\texttt{Path}. In this project, stimulus features are modeled as
$\mathbb{R}^{n}$ time series, EEG/MEG observations are modeled as
$\mathbb{R}^{c \times t}$, source estimates are modeled as
$\mathbb{R}^{m \times n}$, and NCRF coefficients are modeled as
$\mathbb{R}^{m \times i \times \ell}$. Strings are used for identifiers such
as subject, task, stimulus, and estimator name (for example,
\texttt{boosting} and \texttt{ncrf}), while paths denote filesystem locations
for predictors, intermediate artifacts, and outputs. Functions are specified by
explicit input/output type signatures, followed by behavioral requirements and
exception conditions where applicable.

\section{Module Decomposition}

The following table is taken directly from the Module Guide document for this project.

\begin{table}[h!]
\centering
\begin{tabular}{p{0.3\textwidth} p{0.6\textwidth}}
\toprule
\textbf{Level 1} & \textbf{Level 2}\\
\midrule

{Hardware-Hiding Module} & Data Loading Module \\
\midrule

\multirow{3}{0.3\textwidth}{Behaviour-Hiding Module} & Data Preprocessing Module\\
& Estimator Resolution Module\\
& TRF Computation Module\\
\bottomrule

\end{tabular}
\caption{Module Hierarchy}
\label{TblMH}
\end{table}

\newpage

\section{MIS of Hardware-Hiding Module} \label{MIS:HH} 
% \wss{Use labels for cross-referencing}

% \wss{You can reference SRS labels, such as R\ref{R_Inputs}.}

% \wss{It is also possible to use \LaTeX for hyperlinks to external documents.}

\subsection{Module}

M1: Data Loading Module

\subsection{Uses}
OS filesystem APIs, Eelbrain data-loading interfaces.
\subsection{Syntax}

\subsubsection{Exported Constants}
NA

\subsubsection{Exported Access Programs}

\begin{center}
\begin{tabular}{p{3cm} p{4cm} p{4cm} p{3cm}}
\hline
\textbf{Name} & \textbf{In} & \textbf{Out} & \textbf{Exceptions} \\
\hline
\texttt{RawSource.\_load} & \texttt{subject}: \texttt{str}, \texttt{recording}: \texttt{Path}, \texttt{preload}: $\mathbb{B}$ & MNE \texttt{Raw} object & \texttt{FileNotFoundError}, \texttt{OSError} (reader I/O errors) \\
\texttt{save.pickle} & \texttt{res}: Python object, \texttt{dst}: \texttt{Path} & written pickle file at \texttt{dst} & \texttt{OSError}, pickle serialization errors \\
\hline
\end{tabular}
\end{center}

\subsection{Semantics}

\subsubsection{State Variables}

None.

\subsubsection{Environment Variables}

\texttt{fs}: filesystem interface provided by the operating system for reading raw EEG/MEG recordings and writing serialized outputs.

\subsubsection{Assumptions}

\begin{itemize}
\item The operating system provides standard file I/O semantics
(\texttt{Linux}/\texttt{Windows}/\texttt{macOS}).
\item Input recording paths and output destinations are valid and accessible
under current process permissions.
\end{itemize}

\subsubsection{Access Routine Semantics}

\noindent \texttt{RawSource.\_load(subject, recording, preload)}:
\begin{itemize}
\item transition: None.
\item output: Returns an MNE \texttt{Raw} object loaded from the resolved
recording path.
\item exception: Raises \texttt{FileNotFoundError} when the recording cannot be
located, or \texttt{OSError} when low-level reader I/O fails.
\end{itemize}

\noindent \texttt{save.pickle(res, dst)}:
\begin{itemize}
\item transition: Writes serialized object \texttt{res} to destination path
\texttt{dst}.
\item output: Produces a persisted pickle artifact for cache/results reuse.
\item exception: Raises \texttt{OSError} for filesystem write failures and
serialization exceptions for unsupported object graphs.
\end{itemize}

\subsubsection{Local Functions}

None.

\newpage
\section{MIS of Data Preprocessing Module} \label{MIS:Input} 
% \wss{Use labels for cross-referencing}

% \wss{You can reference SRS labels, such as R\ref{R_Inputs}.}

% \wss{It is also possible to use \LaTeX for hyperlinks to external documents.}

\subsection{Module}
MEG/EEG Data Preprocessing and Experiment Configuration.
% \wss{Short name for the module}

\subsection{Uses}
Uses M1 for reading raw MEG/EEG data.


\subsection{Syntax}

\subsubsection{Exported Constants}
% Constants defining predictor file naming conventions.
\begin{itemize}
\item \texttt{PREDICTOR\_FILE\_EXT} = \texttt{".pickle"}
\item \texttt{PREDICTOR\_NAME\_SEP} = \texttt{"\textasciitilde"}
\item \texttt{PREDICTOR\_FILE\_TEMPLATE} = \texttt{"\{stimulus\}\textasciitilde\{predictor\}.pickle"}
\end{itemize}

\subsubsection{Exported Access Programs}

\begin{center}
\begin{tabular}{p{3cm} p{4cm} p{4cm} p{3cm}}
\hline
\textbf{Name} & \textbf{In} & \textbf{Out} & \textbf{Exceptions} \\
\hline
\texttt{load\_epochs} &
\texttt{samplingrate, data, interpolate\_bads} &
Dataset &
\texttt{ValueError} \\

\texttt{add\_predictor} &
dataset, predictor, optional \texttt{filter\_x}, optional target \texttt{y} &
\texttt{ds} extended with predictor column \texttt{code.key} &
\texttt{ValueError}, \texttt{NotImplementedError}, undefined predictor errors \\
\hline
\end{tabular}
\end{center}

\subsection{Semantics}

\subsubsection{State Variables}

\begin{itemize}
\item \texttt{predictors}: map from predictor name (\texttt{str}) to predictor
generator/loader objects.
\end{itemize}

\subsubsection{Environment Variables}

\begin{itemize}
\item Filesystem location \texttt{predictor-dir} containing serialized predictor
files in naming form \texttt{\{stimulus\}\textasciitilde\{predictor\}.pickle}.
\item Optional predictor cache directory for generated predictors.
\item Experiment state keys (for example \texttt{subject}, \texttt{recording},
\texttt{epoch}) used to resolve predictor identity and reproducible randomization.
\end{itemize}

\subsubsection{Assumptions}

\begin{itemize}
\item Input dataset contains event metadata required to identify stimulus
instances (for example task/stimulus label).
\item Predictor definitions are registered in \texttt{predictors} before loading.
\item Time axis requested by TRF routines is valid and compatible with predictor
sampling constraints.
\end{itemize}

\subsubsection{Access Routine Semantics}

\noindent \texttt{add\_predictor(ds, code, filter\_x, y)}:
\begin{itemize}
\item transition: Extends dataset \texttt{ds} by assigning predictor signal(s) to
column \texttt{code.key}. Predictor values are aligned to the response time axis
for each epoch/trial.
\item output: None (in-place update of \texttt{ds}).
\item exception:
\begin{itemize}
\item Undefined predictor code (not registered in \texttt{predictors}).
\item \texttt{ValueError} when \texttt{filter\_x} is requested for unsupported
predictor types.
\item \texttt{NotImplementedError} for unsupported combinations (for example
event predictor with variable-time epochs, or session predictor with fixed-time
epochs).
\end{itemize}
\end{itemize}

\noindent \texttt{load\_epochs(samplingrate, data, interpolate\_bads)}:
\begin{itemize}
\item transition: No module state transition. Reads epoch definitions and raw
recordings to construct an analysis dataset for the current subject/recording.
\item output: Returns epoch dataset \texttt{ds} containing response variable
(for example MEG/EEG channels), event metadata, and a time axis consistent with
\texttt{samplingrate}.
\item exception:
\begin{itemize}
\item \texttt{FileNotFoundError}/\texttt{OSError} when raw/event files are
missing or unreadable.
\item \texttt{ValueError} for invalid data-mode or inconsistent epoch
configuration.
\item Propagates preprocessing/runtime errors from lower-level epoch-loading
utilities.
\end{itemize}
\end{itemize}

\subsubsection{Local Functions}

\begin{itemize}
\item \texttt{load\_predictor(code)}: parse predictor expression and normalize
stimulus/shuffle qualifiers.

\end{itemize}

\newpage
\section{MIS of Estimator Resolution Module} \label{MIS:Estimator}

\subsection{Module}

Estimator resolution service for selecting estimator policies and producing
effective TRF options that determine downstream data flow.
\subsection{Uses}
Uses M2 for experiment state/configuration and M4 for TRF fitting execution.

\subsection{Syntax}

\subsubsection{Exported Constants}
\begin{itemize}
\item \texttt{ESTIMATOR\_KEYS} = \{\texttt{"boosting"}, \texttt{"ncrf"}\}
\end{itemize}

\subsubsection{Exported Access Programs}

\begin{center}
\begin{tabular}{p{3cm} p{4cm} p{4cm} p{3cm}}
\hline
\textbf{Name} & \textbf{In} & \textbf{Out} & \textbf{Exceptions} \\
\hline
\texttt{\_resolve\_estimator} &
\texttt{estimator: str|Estimator|None} &
\texttt{Estimator|None} &
\texttt{ValueError} (invalid estimator key or type) \\
\hline
\end{tabular}
\end{center}

\subsection{Semantics}

\subsubsection{State Variables}

None. M3 only specifies the in-function resolution behavior.

\subsubsection{Design Pattern}

Strategy pattern with a registry-based resolver: the estimator key selects a
concrete estimator object whose parameters are applied before TRF computation.

\subsubsection{Environment Variables}

\begin{itemize}
\item Filesystem cache for TRF artifacts (read/write via path resolution).
\item Forward model/covariance inputs for NCRF path when source-space inversion
is requested.
\end{itemize}

\subsubsection{Assumptions}

\begin{itemize}
\item If \texttt{estimator} is provided as a string, the experiment defines a
dictionary \texttt{estimators} and contains the requested key.
\item Estimator objects implement \texttt{parameters\_for\_partial()} and return
keyword arguments compatible with the selected backend fitter.
\item Predictor and response data are already prepared by upstream modules.
\item Estimator resolution occurs before the TRF fitting job is created.
\end{itemize}

\subsubsection{Access Routine Semantics}

\noindent \texttt{\_resolve\_estimator(estimator)}:
\begin{itemize}
\item transition: None.
\item output: Returns a resolved estimator instance (or \texttt{None}) and
validates estimator type and key.
\item exception:
\begin{itemize}
\item \texttt{ValueError} if \texttt{estimator} is a string but
\texttt{estimators} is missing/not a dictionary.
\item \texttt{ValueError} if estimator key is not found in
\texttt{estimators} or if the resolved object is not an \texttt{Estimator}.
\end{itemize}
\end{itemize}

\subsubsection{Local Functions}

\begin{itemize}
\item \texttt{\_apply\_estimator\_params(estimator: Estimator|None, **kwargs: dict)}:
merge estimator defaults into explicit TRF options.
\end{itemize}

\newpage

\section{MIS of TRF Computation Module} \label{MIS:TRF}

\subsection{Module}

TRF computation and backend dispatch.
\subsection{Uses}
Uses M2 for epoch/predictor-formatted data and M3 for resolved estimator
options.


\subsection{Syntax}

\subsubsection{Exported Constants}
NA

\subsubsection{Exported Access Programs}

\begin{center}
\begin{tabular}{p{3cm} p{4cm} p{4cm} p{3cm}}
\hline
\textbf{Name} & \textbf{In} & \textbf{Out} & \textbf{Exceptions} \\
\hline
\texttt{load\_trf} & 
\texttt{x}, \texttt{tstart}, \texttt{tstop}, TRF options, optional \texttt{make}, optional \texttt{path\_only} &
TRF result object, or cache path if \texttt{path\_only=True} &
\texttt{IOError} (cache miss with \texttt{make=False}), propagated cache/load/runtime errors \\
\texttt{\_trf\_job} & 
effective TRF options, prepared response \texttt{y} and predictor(s) \texttt{x} &
A deferred callable created with \texttt{partial(...)}\, representing a TRF fitting task (e.g., \texttt{boosting} or NCRF). &
\texttt{ValueError} (invalid backward multi-predictor case), runtime fitting exceptions from backend libraries \\
\hline
\end{tabular}
\end{center}

\subsection{Semantics}

\subsubsection{State Variables}

None. Computation is request-scoped.

\subsubsection{Environment Variables}
\begin{itemize}
\item Filesystem cache for TRF artifacts (read/write via path resolution).

\end{itemize}

\subsubsection{Assumptions}

\begin{itemize}
\item Response epochs and predictors requested by model \texttt{x} are available
from upstream modules.
\item Effective estimator options have already been resolved (M3) before
dispatching the fit callable.
\item Cache keys uniquely encode effective TRF options.
\end{itemize}

\subsubsection{Access Routine Semantics}

\noindent \texttt{load\_trf(x, tstart, tstop, ..., estimator, make, path\_only)}:
\begin{itemize}
\item transition: May create or update cache file at resolved \texttt{trf-file}
path when \texttt{make=True}; otherwise reads existing cache only.
\item output:
\begin{itemize}
\item returns cache path if \texttt{path\_only=True};
\item returns cached TRF if valid cache exists;
\item otherwise calls \texttt{\_trf\_job(...)} to obtain a deferred fit callable
and executes it (\texttt{res = func()}), then serializes result to cache.
\end{itemize}
\item exception:
\begin{itemize}
\item \texttt{IOError} if requested TRF is absent and \texttt{make=False}.
\item Propagates unpickling and backend runtime errors.
\end{itemize}
\end{itemize}

\noindent \texttt{\_trf\_job(x, tstart, tstop, ..., estimator)}:
\begin{itemize}
\item transition: None in module state.
\item output: Returns deferred computation using \texttt{partial(...)}:
\begin{itemize}
\item boosting path:
\[
\texttt{partial(boosting, y, xs, tstart, tstop, 'inplace', **boosting\_args)}
\]
\item NCRF path:
\[
\texttt{partial(fit\_ncrf, y, xs, fwd, cov, tstart, tstop, normalize=True, in\_place=True, **ncrf\_args)}
\]
Estimator-specific options are merged into fitter kwargs before creating the
\texttt{partial} callable.
\item \texttt{None} when equivalent NCRF output already exists and cache link is
created.
\end{itemize}
\item exception:
\begin{itemize}
\item \texttt{ValueError} for backward model with multiple predictors.
\item Propagates fitting/library runtime errors from backend calls.
\end{itemize}
\end{itemize}

\subsubsection{Local Functions}

\begin{itemize}
\item \texttt{build\_boosting\_partial(y, xs, tstart, tstop, boosting\_args)}:
constructs deferred boosting computation.
\item \texttt{build\_ncrf\_partial(y, xs, fwd, cov, tstart, tstop, ncrf\_args)}:
constructs deferred NCRF computation.
\end{itemize}

\newpage

\bibliographystyle {plainnat}
\bibliography {../../../refs/References}

\newpage



% \section{Appendix} \label{Appendix}

% \wss{Extra information if required}

% \newpage{}

% \section*{Appendix --- Reflection}

% \wss{Not required for CAS 741 projects}

% The information in this section will be used to evaluate the team members on the
% graduate attribute of Problem Analysis and Design.

% \input{../../Reflection.text}

% \begin{enumerate}
%   \item What went well while writing this deliverable? 
%   \item What pain points did you experience during this deliverable, and how
%     did you resolve them?
%   \item Which of your design decisions stemmed from speaking to your client(s)
%   or a proxy (e.g. your peers, stakeholders, potential users)? For those that
%   were not, why, and where did they come from?
%   \item While creating the design doc, what parts of your other documents (e.g.
%   requirements, hazard analysis, etc), it any, needed to be changed, and why?
%   \item What are the limitations of your solution?  Put another way, given
%   unlimited resources, what could you do to make the project better? (LO\_ProbSolutions)
%   \item Give a brief overview of other design solutions you considered.  What
%   are the benefits and tradeoffs of those other designs compared with the chosen
%   design?  From all the potential options, why did you select the documented design?
%   (LO\_Explores)
%   \item (After you have implemented another team's module, which means this
%   isn't filled in until after the original deadline). What did you learn by
%   implementing another team's module? Were all the details you needed in the
%   documentation, or did you need to make assumptions, or ask the other team
%   questions?  If your team also had another team implement one of your modules,
%   what was this experience like?  Are there things in your documentation you
%   could have changed to make the process go more smoothly for when an
%   ``outsider'' completes some of the implementation?
% \end{enumerate}


\end{document}
